% STATUS: written 2026-11-06 (English gate clean; sources read in full as recorded below).
% Sources read in full: Jacobson 1937 (Trans. AMS 42, 206--224).
\section[Derivations and Lie algebras: Jacobson 1937 and the structure theory][Derivations and Lie algebras]{Derivations and Lie algebras: Jacobson 1937 and the structure theory}\label{sec:origins:lie-theory}

The first systematic study of derivations of an abstract algebra is a paper
of Jacobson from 1937~\cite{Jacobson1937}. Its first sentence states its
purpose. It studies ``the algebraic properties of the set of operations
mapping an algebra on itself and having the formal character of derivation
in the field of analytic functions''~\cite[p.~206]{Jacobson1937}.

Jacobson compares his results with the theory of automorphisms. He writes:
``Some of the results obtained are analogous to well-known theorems on
automorphisms of algebras''~\cite[p.~206]{Jacobson1937}. In analysis, a
derivation is an automorphism taken at an infinitely small scale. His own
footnote gives an example. On $k[x]$, the map $\exp D$ acts as
$f(x) \mapsto f(x+1)$~\cite[p.~206]{Jacobson1937}.

Jacobson asks two questions. What do the derivations of an abstract algebra
look like? Which theorems on automorphisms also hold for derivations? He has
no analysis to help him. There are no limits and no convergence, and the
ground field is arbitrary.

Two facts about derivations have no match for automorphisms. They drive the
paper. First, derivations are closed under commutators, so they form a Lie
algebra rather than a group. Second, in characteristic $p$ they are also
closed under $p$-th powers.

To answer these questions, Jacobson works with any algebra $A$ over a field $k$. The algebra need not be
commutative. It need not be associative. It need not be finite dimensional.

\begin{definition}[{\cite[p.~206]{Jacobson1937}}]\label{def:origins:jacobson}
A \emph{derivation} of $A$ is a $k$-linear map $D\colon A \to A$ with
\[
D(xy) = D(x)\,y + x\,D(y) \qquad (x, y \in A).
\]
\end{definition}

Jacobson wrote maps on the right, as $(xy)D = (xD)y + x(yD)$. This was usual
in the algebra of his time. In this book we write maps on the left.

The first properties follow at once from the product rule. The sum of two derivations is a derivation. A scalar multiple of a
derivation is a derivation. The product $D_1D_2$ of two derivations is
usually not a derivation. But the commutator $[D_1,D_2] = D_1D_2 - D_2D_1$ is
a derivation~\cite[p.~207]{Jacobson1937}.

Repeated use of the product rule gives the Leibniz formula:
\[
D^n(xy) = \sum_{i=0}^{n} \binom{n}{i} D^{i}(x)\, D^{n-i}(y).
\]
Now let the characteristic be $p$. Then the coefficients $\binom{p}{i}$ with
$0 < i < p$ are zero. So $D^p$ is again a
derivation~\cite[p.~207, (7)--(8)]{Jacobson1937}. This fact leads to a new
notion. A \emph{restricted Lie algebra} is a space of linear maps with two
properties. It is closed under commutators. In characteristic $p$, it is also
closed under $p$-th powers.

\begin{theorem}[{\cite[Theorem~1, p.~207]{Jacobson1937}}]\label{thm:origins:jacobson-restricted}
The derivations of an algebra $A$ over $k$ form a restricted Lie algebra
$\operatorname{Der}(A)$ of linear maps of~$A$.
\end{theorem}

Jacobson calls $\operatorname{Der}(A)$ the derivation algebra of $A$.

Jacobson also returns to the comparison with automorphisms. The link is the
exponential. In a footnote he takes the polynomial ring $k[x]$ over a field
of characteristic $0$, with the usual derivative $D(f) = f'$. Then
$\exp D = 1 + D + D^2/2! + \cdots$ is an automorphism, since
$(\exp D)(f)(x) = f(x+1)$~\cite[p.~206, footnote]{Jacobson1937}.
Section~\ref{sec:basic-theory:exponential} treats this link in general.

The rest of the paper turns to associative algebras of finite order, mainly
semisimple ones. Here Jacobson uses the known theory of algebras. He refers
the reader to the book \emph{Algebren} of Deuring~\cite[p.~206]{Jacobson1937}.

A recent line of work pairs an algebra with one chosen derivation and studies
the pair. Das calls such a pair a \emph{LeibDer pair} when the algebra is a
Leibniz algebra and the map is a derivation of
it~\cite[\S 3, p.~4]{Das2021LeibDer}. With a skew-symmetric bracket the same
construction gives a LieDer pair, so a LeibDer pair is the version without
skew-symmetry~\cite[\S 3, p.~4]{Das2021LeibDer}.

\begin{definition}[{\cite[\S 3, p.~4]{Das2021LeibDer}}]\label{def:origins:leibder}
Let $\mathfrak{g}$ be a Leibniz algebra. A linear map
$\varphi_{\mathfrak{g}}\colon \mathfrak{g} \to \mathfrak{g}$ is a
\emph{derivation} on $\mathfrak{g}$ when
\[
  \varphi_{\mathfrak{g}}[x, y] = [\varphi_{\mathfrak{g}}(x), y]
  + [x, \varphi_{\mathfrak{g}}(y)]
\]
for all $x, y \in \mathfrak{g}$. A \emph{LeibDer pair}
$(\mathfrak{g}, \varphi_{\mathfrak{g}})$ is a Leibniz algebra together with
one of its derivations.
\end{definition}

Das notes that the derivations in this sense are the one-cocycles of the
cochain complex of the Leibniz algebra with coefficients in
itself~\cite[\S 3, p.~4]{Das2021LeibDer}. From there the pair carries its own
cohomology, and central extensions, deformations and the homotopy version are
all read through it~\cite[p.~1]{Das2021LeibDer}. For a free Leibniz algebra on
a space $V$, a linear map of $V$ extends to a derivation, and the resulting
pair is the free one~\cite[Proposition~3.4, p.~5]{Das2021LeibDer}.

The derivations of Jacobson map an algebra to itself. Derivations with values
in a module have their own history. It began with groups that act on fields,
long before 1937.
